\begin{figure} [h]
	\centering
	
	\tikzstyle{comp} = [circle, draw, minimum size = 0.5cm, line width=0.7] % comparator style
	\tikzstyle{block} = [rectangle, draw, minimum width = 1.2cm, minimum height = 1.2cm, line width=0.7]  % bloc style
	\tikzstyle{big_block} = [rectangle, draw, minimum width = 2cm, minimum height = 2cm, line width=0.7]  % bloc style
	
	\newcommand{\comp}[3][]{ % comparator command
		\node[comp] (#2) at #3 {};
		\draw (#2) -- + (45:0.25) -- + (45:-0.25);
		\draw (#2) -- + (-45:0.25) -- + (-45:-0.25);
		\ifthenelse{
			\isempty{#1}{}
		}{}
		{
			\ifthenelse{
				\equal{#1}{a}
			}{
				\node[above right] at (#2) {$-$\phantom{p}};
			}{}
			\ifthenelse{
				\equal{#1}{b}
			}{
				\node[below right] at (#2) {$-$\phantom{l}};
			}{}
		}
	}
	\tikzsetnextfilename{fig_impedance_simulation_real_position}		
	\begin{tikzpicture} [node distance = 2cm]	
		
		\comp[]{c1}{(0, 0)};
		
		\node[block, left of = c1] (imp_model_1) {$\frac{1}{K + Bs + Ms^2}$};
				
		\draw[-{Latex[length=3mm]}] (imp_model_1.east) -- (c1.west);
		
		\draw[-{Latex[length=3mm]}] (c1.east) --+ (1,0)
			node[align = center, below, midway] {$\vect{x}_{z}$};
			
		\draw[{Latex[length=3mm]}-] (imp_model_1.west) --+ (-1,0)
			node[align = center, below, midway] {$\vect{f}_{p}$};
			
		\draw[{Latex[length=3mm]}-] (c1.south) --+ (0,-1) --+ (-4,-1) 
			node[align = center, below] {$\vect{x}_{z0}$};	

	\end{tikzpicture}
	\caption{Simulation pour l'identificaiton des paramètres en impédance avec des signaux en force issus d'expérimentations sans perturbations.}
	\label{fig_impedance_simulation_real_position} % Label 
\end{figure}